\section{R ile çözümlü uygulama}
\label{sec:r-b08}

\textbf{Soru.} $n=25$, $\bar x=72$ ve $s=10$ için evren ortalamasının
yüzde 95 ve yüzde 99 güven aralıklarını hesaplayın. Burada ham gözlemler
değil özet istatistikler verildiği için $t$ kritik değeriyle doğrudan hesap yapın.

\begin{lstlisting}[style=prismR]
hacim <- 25
ortalama <- 72
standart_sapma <- 10
standart_hata <- standart_sapma / sqrt(hacim)
serbestlik <- hacim - 1
guven_araligi <- function(duzey) {
  kritik <- qt((1 + duzey) / 2, df = serbestlik)
  ortalama + c(-1, 1) * kritik * standart_hata
}
c(standart_hata = standart_hata,
  kritik_95 = qt(0.975, df = serbestlik))
rbind(GA95 = guven_araligi(0.95),
      GA99 = guven_araligi(0.99))
\end{lstlisting}

\begin{outputguidebox}
\textbf{Çözüm ve kontrol değerleri.} Standart hata $10/\sqrt{25}=2$,
serbestlik derecesi 24 ve yüzde 95 aralık için kritik değer 2,0639'dur.
Yüzde 95 güven aralığı $[67{,}8722;76{,}1278]$, yüzde 99 güven aralığı
yaklaşık $[66{,}4061;77{,}5939]$ olur.
\end{outputguidebox}

\textbf{Yorum.} Bağımsız gözlemler ve küçük örneklemde yaklaşık normal
bir evren modeli altında bu $t$ aralığı kullanılır. Özet istatistikler aykırı
değerleri ve dağılım biçimini denetlemeye yetmez. Ham veri elde olsaydı
\texttt{t.test(puan, conf.level = 0.95)\$conf.int} kullanılabilirdi;
\texttt{puan} yerine üç özet sayıyı koymak yanlış olurdu.

\textbf{Pekiştirme ve yanıt.} Güven düzeyi arttıkça kritik değer ve aralık
genişliği artar. Yüzde 95 ifadesi, sabit parametrenin hesaplanmış bu aralıkta
bulunma olasılığı değil yöntemin uzun dönem kapsama oranıdır.